\documentclass[10pt,a4paper]{amsart}
\usepackage[margin=1in]{geometry}
\usepackage[T1]{fontenc}
\usepackage{lmodern,microtype,amsmath,amssymb,graphicx,booktabs,array,tikz,flafter,needspace}
\usepackage[colorlinks=true,linkcolor=blue!45!black,citecolor=blue!45!black,urlcolor=blue!45!black]{hyperref}
\usetikzlibrary{arrows.meta,decorations.markings}
\newtheorem{theorem}{Theorem}
\newtheorem{proposition}[theorem]{Proposition}
\newtheorem{lemma}[theorem]{Lemma}
\theoremstyle{remark}
\newtheorem*{example}{Example}
\newcommand{\eps}{\varepsilon}
\newcommand{\In}{\operatorname{In}}
\newcommand{\cw}{\mathrm{cw}}
\newcommand{\ccw}{\mathrm{ccw}}
\newcommand{\fig}[3]{\begin{figure}[htbp]\centering\includegraphics[trim=1 0 1 0,clip,width=#2\textwidth]{../figures/#1.pdf}\caption{#3}\label{fig:#1}\end{figure}}
\setlength{\emergencystretch}{2em}
\title{Hamiltonian cycles from states}
\author{An exposition of OpenAI's proof}
\date{November 14, 2026}
\begin{document}
\begin{abstract}
A spanning subgraph of degree two need not be a single cycle.
We explain OpenAI's proof of Barnette's conjecture by constructing a forest in the planar dual that forces the complementary degree-two subgraph to be connected.
Two disjoint face--vertex matchings encode the local choices.
A signed count cancels every cyclic choice, while heights at white faces show that the count is nonzero.
The surviving choice yields two induced trees with a prescribed facial pattern, which allows the construction to pass across separating triangles.
\end{abstract}
\maketitle

\section{Degree two and connectivity}
A Hamiltonian cycle is a connected spanning subgraph of degree two. In a cubic graph, choosing the omitted edge at each vertex is the easy part: the complement of a perfect matching has degree two everywhere. Connectivity is the whole problem. On the cube, for instance, the two opposite square faces give a spanning subgraph of degree two with two components.

We will enforce connectivity in the dual graph. There, omitted edges become a forest. Planarity turns acyclicity of those edges into connectivity of their complement. The local degree condition and the global connectivity condition then have separate certificates.

All graphs in the theorem are finite, simple and undirected. A graph is cubic if every vertex has degree three, and is 3-vertex-connected if it has at least four vertices and remains connected after deleting any two or fewer vertices.
\begin{theorem}[OpenAI \cite{openai2026}]\label{thm:main}
Let $G$ be a cubic, bipartite, planar, 3-vertex-connected graph. For every edge $e$ of $G$, there is a Hamiltonian cycle avoiding $e$.
\end{theorem}
The edge prescription will let solutions on smaller pieces agree when we glue them. The proof below explains OpenAI's paired-state argument. Barnette's conjecture was recorded in the 1969 Waterloo proceedings \cite{barnette1969}. Stein's region-partition formulation \cite[Theorem 2.2]{stein1971} led to the two-induced-tree formulation given by Alt, Payne, Schmidt and Wood \cite[Lemma 5]{apsw2016}. The prescribed-edge strengthening and the reduction to graphs without nontrivial three-edge cuts belong to the earlier work of Kelmans, as explained by Hertel and by Alt et al.\ \cite{hertel2005,apsw2016}.

Our main change is in the signed count. Li \cite{li2026} replaces OpenAI's exponential sum by involutions on marked crossings. The constructive exposition \cite{constructive2026} uses integer signs and a polynomial. The root-component count below also follows that exposition. We use the product of face heights proposed in the project note \cite[Sections 1--3]{heightnote}, and prove its counting identity here. We also use that note's formulation of the disk identity and its short recovery argument. The product of heights uses the same cancellation mechanism, with each factor detecting enclosure of one face.

Fix a sphere embedding of $G$, drawn with piecewise linear edges, and let $T=G^*$ be its dual. This is a simple triangulation. Here are the details needed for simplicity. A cut of one or two edges in $G$ would have at least three vertices on either side: a singleton in a simple cubic graph has three outgoing edges and a pair has at least four. Deleting the at most two cut endpoints on one side would therefore disconnect a nonempty remainder, contrary to 3-connectivity. A loop or a pair of parallel edges in the dual would draw a Jordan curve crossing just one or two primal edges, and would give such a cut. Cubicity now makes every dual facial boundary a closed walk of length three in a simple graph, hence a triangle. Euler's formula gives $|V(T)|=2+|V(G)|/2\geq4$. The bipartition of $G$ colors the faces of $T$ black and white, with opposite colors across every edge.

We use the following elementary fact about duality. A spanning forest includes all vertices, with isolated vertices allowed.
\Needspace{8\baselineskip}
\begin{lemma}\label{lem:duality}
In a connected sphere graph, the edges complementary to a spanning tree dualize to a spanning tree. The edges complementary to a spanning forest dualize to a connected spanning subgraph.
\end{lemma}
\begin{proof}
Contract the tree to one vertex. Faces are preserved, and its edges disappear from the dual. The remaining dual is connected and has $|E|-|V|+1=|F|-1$ edges by Euler's formula, so it is a tree. Extend a forest to a spanning tree to obtain the second assertion.
\end{proof}
Thus we want a forest $P$ in $T$ with exactly one side of each triangular face. Its complementary dual has degree two because each triangle has two sides left, and is connected by Lemma~\ref{lem:duality}.

For gluing, we will obtain this forest as two induced trees. For $U\subseteq V(T)$, the induced graph $T[U]$ contains every edge whose endpoints lie in $U$; a single vertex counts as a tree. If two induced trees partition $V(T)$, no triangle is monochromatic. Each triangle therefore has exactly one edge within a class, as required. We prove the following precise dual statement.
\Needspace{7\baselineskip}
\begin{proposition}\label{prop:dual}
Let $T$ be a simple sphere triangulation with at least four vertices and a proper black-and-white face coloring. Given any facial triangle $t_0$ and any side $e_0$ of it, the vertices partition into two induced trees such that the endpoints of $e_0$ belong to one tree and the third corner of $t_0$ belongs to the other.
\end{proposition}
We first construct the partition when there is no separating triangle, meaning a 3-cycle with vertices strictly on both sides. Section~\ref{sec:glue} then proves the proposition by induction. The matching and cancellation arguments in between will also apply whenever a pair of states exists.

\begin{example}[The cube]
Its dual is the octahedron. Use roots $0,1,2$, other vertices $3,4,5$, and the following faces, written here as unordered vertex sets:
\[
\begin{array}{c|cccc}
\text{black}&t_0=012&t_1=135&t_2=234&t_3=045\\
\text{white}&y_0=015&y_1=123&y_2=024&y_3=345.
\end{array}
\]
We prescribe $e_0=01$, so the eventual cube cycle must omit $t_0y_0$. The octahedron labels in Figure~\ref{fig:01-cube} are these vertex labels. We will follow this example from states to the recovered cycle.
\end{example}
\fig{01-cube}{1}{The cube construction, read from left to right along each row. The first panel shows two cycles rather than one. In the octahedron the roots are $0,1,2$. Outlined and filled wedges mark the two states; the opposite edges form three rooted trees. Deleting the filled spokes gives an incidence tree and a complementary tree on white faces. The last panel shows the selected forest and the single dual cycle. Later sections justify each transition.}

\section{Local choices with a second state}\label{sec:states}
Make $t_0$ the outer face, and exchange the colors if needed so that it is black. Its three corners are the roots $R$. If there are $k$ black faces, side counting gives $k$ white faces and $3k$ edges. Euler's formula gives $k+2$ vertices. Consequently the bounded black faces $A$ and the nonroot vertices $X$ have the same size,
\begin{equation}\label{eq:counts}
 |A|=|X|=n=k-1,\qquad |V(T)|=n+3.
\end{equation}

We need a local choice whose acyclicity can eventually be tested. Put a center $c_t$ in each black face and join it to its corners by three spokes. We will delete one spoke at each bounded black center and one at each nonroot. Such deletions are a perfect matching in the bipartite incidence graph $B$ on $A\cup X$, where $t$ and $v$ are adjacent when $v$ is a corner of $t$. Equivalently, a \emph{state} is a bijection $s:A\to X$ with $s(t)\in V(t)$. These are Tutte states \cite{tutte1975,hine-kalman2018}. We prove the state properties used below.

After deleting the spoke to $s(t)$, the two remaining spokes join the other corners. Suppress the center and replace that path by the opposite side of the triangle. Call the resulting spanning graph $Q_s$. It has one edge per face of $A$. Distinct black faces supply distinct edges, since an edge has only one black incident face.

A forest $Q_s$ would be useful only if its components keep the three roots separate. A second state supplies exactly this control. A \emph{pair} $(r,s)$ consists of states with $r(t)\ne s(t)$ at every $t$. If $u(t)$ is the third corner, orient the edge of $Q_s$ from $r(t)$ to $u(t)$. Every nonroot then has outdegree one, and every root has outdegree zero.
\Needspace{8\baselineskip}
\begin{lemma}\label{lem:components}
If $s$ has a partner $r$, each component of $Q_s$ is either a tree with one root or a graph with a unique cycle and no root. If $Q_s$ is a forest, it has three components, one at each root, and $s$ has exactly one partner.
\end{lemma}
\begin{proof}
In a component with $v$ vertices and $a$ roots, summing outdegrees gives $v-a$ edges. Connectedness gives at least $v-1$ edges, so $a\leq1$. For $a=1$ the component is a tree; for $a=0$ it has one cycle. In a rooted tree all edges must point towards its root. In a rootless component all off-cycle edges point towards the cycle, and the cycle has either of its two directions. These orientations determine the partner by taking the tail in each black triangle. In particular the forest case gives a unique partner.
\end{proof}
The same outdegree observation will be used repeatedly: in a graph with outdegree at most one, an undirected simple cycle must be directed. Its edges require distinct tails, exhausting all its vertices.

\begin{example}[The cube continued]
In the face order $t_1,t_2,t_3$ above, take $r=(3,4,5)$ and $s=(5,3,4)$. The opposite edges of $s$ are $13,24,05$. They give three rooted trees. Both states of this rooted octahedron have forest opposite-edge graphs. To see cancellation of a cyclic $Q_s$, we need a larger example.
\end{example}
For that purpose we also use the 13-vertex triangulation in Figure~\ref{fig:03-dual}, whose dual has 22 vertices. We retain the figure's vertex and face labels throughout. Its black faces are $t_0=054$ and the ten columns below. The symbols $a_j$ are the numerical face identifiers used in the figures.
\begin{equation}\label{eq:running}
\begin{array}{c|rrrrrrrrrr}
 &a_0&a_2&a_7&a_9&a_{11}&a_{13}&a_{15}&a_{17}&a_{19}&a_{21}\\
 V(a_j)&016&032&187&129&349&4\,11\,10&56\,11&67\,12&8\,11\,12&89\,10\\
 r(a_j)&6&2&8&1&3&10&11&7&12&9\\
 s(a_j)&1&3&7&2&9&11&6&12&8&10
\end{array}
\end{equation}
Here $R=\{0,4,5\}$ and $e_0=04$. The edges $18,19,89$ form a cycle in $Q_s$, whereas $Q_r$ is a forest. These assignments and all examples below can be checked directly from the face lists; the companion calculation \cite{draftdata} also recomputes them with the reference implementation.
\fig{03-dual}{.72}{The larger running example. The dark triangulation has roots $0,4,5$ and ten nonroots. Shaded triangles are black, including the outer face indicated by the collar. Each grey vertex of $G$ lies in one triangle and each grey edge crosses its corresponding edge of $T$. The same labels are used in~\eqref{eq:running} and in the recovery figures.}

\fig{09-Q}{1}{The two opposite-edge graphs from~\eqref{eq:running}. The dashed edges on the left contain the triangle $1,8,9$, with no root in its component. On the right, $Q_r$ has three tree components, one at each root. The proof must guarantee a pair with the second behavior.}

We now have two tasks: obtain a pair at all, then find a pair whose second state has forest $Q_s$. The following density argument is OpenAI's \cite[Proposition 3.1]{openai2026}.
\Needspace{8\baselineskip}
\begin{lemma}\label{lem:pair}
If $T$ has no separating triangle, then a pair exists.
\end{lemma}
\begin{proof}
Give capacity two to edges from a source to $X$ and from $A$ to a sink, and capacity one to each incidence from $X$ to $A$. For a cut with $S\subseteq X$ on the source side, minimizing over the placement of the black faces gives capacity
$2(n-|S|)+\sum_{t\in A}\min(2,|V(t)\cap S|)$.
Thus it suffices to prove
\begin{equation}\label{eq:flow}
 \sum_{t\in A}\min(2,|V(t)\cap S|)\geq2|S|\qquad(S\subseteq X).
\end{equation}
Integral max-flow/min-cut \cite{ford1956} then gives a flow of value $2n$. All source and sink edges are saturated. The incidences carrying one unit form a spanning degree-two subgraph of $B$, whose even cycles split alternately into two disjoint perfect matchings. For completeness, integrality follows by augmenting from zero along residual paths: the bottleneck is an integer, the value increases each time, and if no path remains the reachable vertices give a cut of equal value.

To prove~\eqref{eq:flow}, put $U=V(T)\setminus S$. It contains $R$. Write $e(U)$ for the number of edges induced by $U$ and $b(U)$ for the number of actual black faces all of whose corners lie in $U$. Include $t_0$ in the sum with contribution zero. A black face with $j$ corners in $U$ has shortfall $\max(0,j-1)$ from two. This is its number of induced sides minus one when $j=3$. Summing over black faces and using~\eqref{eq:counts} shows that~\eqref{eq:flow} is equivalent to
\begin{equation}\label{eq:density}
 e(U)-b(U)\leq2|U|-4.
\end{equation}

Consider a connected component $K$ of $T[U]$ with $p\geq4$ vertices, $m$ edges and $f$ faces in its sphere embedding. Face lengths count sides, so a bridge counts twice. Simplicity and connectedness imply that all face lengths are at least three. Every triangular face of $K$ is an actual face of $T$: its boundary is a nonseparating 3-cycle, one side is an actual face of $T$, and $K$ has a vertex off that cycle on the other side. Conversely an actual face with its corners in $K$ is a face of $K$.

Let $b,w$ count these black and white triangular faces. In any other face region of $K$, of length $d\geq4$, let $b',w'$ count the original faces of $T$ in that region. Counting signed sides gives $3(b'-w')=\sum_{j=1}^d\sigma_j$, with each $\sigma_j\in\{1,-1\}$. Internal sides cancel; this remains true with repeated boundary sides and with other components inside the region. Therefore $|b'-w'|\leq d-4$: for $d=4$ parity and divisibility force the sum to be zero, for $d=5$ its absolute value is at most three, and for $d\geq6$ use $d/3\leq d-4$.

The regions partition the original faces, whose total color balance is zero. Hence $w-b=\sum(b'-w')\leq\sum(d-4)$, with sums over the nontriangular regions. Counting sides and applying Euler gives
\[
 2m=4f-b-w+\sum(d-4)\geq4f-2b,
 \qquad p-m+f=2,
 \qquad m-b\leq2p-4.
\]
The component containing $R$ also satisfies this last bound if it has only three vertices: it is the triangle $t_0$, with three edges and one actual black face. Every other component satisfies the weaker bound $m-b\leq2p$, by the preceding argument for $p\geq4$ and by simplicity for $p\leq3$. Summing proves~\eqref{eq:density}.
\end{proof}
\section{What a state cycle encloses}\label{sec:disk}
The second task is global: cycles in $Q_s$ must disappear. We will cancel pairs with cyclic $Q_s$ in a signed count. To choose the sign, we first need a restriction on cycles made from state choices. It will also prove acyclicity at the end of the construction.

For an oriented edge $v\to w$ of $T$, let $\delta(v,w)=1$ when its black face lies on the left, and $-1$ otherwise. Thus $\delta(w,v)=-\delta(v,w)$.
\Needspace{8\baselineskip}
\begin{lemma}[Disk identity]\label{lem:disk}
Let $r$ be a state, and choose in each $t\in A$ an edge directed from $r(t)$ to another corner. A simple directed cycle $C$ of chosen edges bounds a disk containing exactly one more white face than black faces. Write $B,W$ for the numbers of black and white faces in that disk. If its length is $h$ and $\ell$ boundary edges have their black face inside, then
\begin{equation}\label{eq:disk}
 W=B+1,\qquad h-2\ell=3.
\end{equation}
In particular $h$ is odd. The $\delta$-sum along $C$ is $-3$ counterclockwise and $3$ clockwise.
\end{lemma}
\begin{proof}
Only nonroots have outgoing edges, so $C$ avoids all roots. Its bounded disk contains no root, since each root is incident with the unbounded face $t_0$. Write $I$ for the number of strictly interior vertices and $E_d$ for the number of edges in the closed disk, including chords. Here $B,W$ are face counts, not the incidence graph $B$.

The state matches each black face inside to an interior or boundary vertex. Every interior vertex is matched to a face inside. A boundary vertex is matched to the black face of its outgoing cycle edge, which lies inside in exactly $\ell$ cases. Bijectivity gives the first of the three counts
\begin{equation}\label{eq:threecounts}
 B=I+\ell,\qquad E_d=3B+h-\ell=3W+\ell,
 \qquad (I+h)-E_d+B+W=1.
\end{equation}
For the middle count, each edge has one black and one white incident face; count the edges by the black faces inside, then add the boundary edges with black outside. The last count is Euler's formula for the disk. Substitution gives $W-B=1$, and the two expressions for $E_d$ give $h-2\ell=3$. Counterclockwise the disk is on the left, so the $\delta$-sum is $\ell-(h-\ell)=-3$; reversing the direction changes its sign.
\end{proof}
This is OpenAI's disk identity \cite[Lemma 4.1]{openai2026}, expressed as a face excess. The matching count $B=I+\ell$ is the reason it holds for state cycles. Euler's formula and face coloring alone do not give the conclusion.

\begin{example}[The larger disk]
Join $r(a_j)$ to $s(a_j)$ in~\eqref{eq:running}. One directed cycle is $1,2,3,9,10,11,6,1$, counterclockwise; the other is $8,7,12,8$, clockwise. The long cycle encloses $B=5$ black faces, $W=6$ white faces and $I=3$ vertices, namely $7,8,12$. It has $h=7$ and $\ell=2$, so all three counts in~\eqref{eq:threecounts} can be seen in Figure~\ref{fig:08-disk}. The smaller cycle encloses the white triangle $7,8,12$.
\end{example}
\fig{08-disk}{.85}{The long cycle in the running example encloses five black and six white faces. Three matched vertices lie strictly inside; two more are tails of boundary edges whose black face lies inside. Thus the state contributes $B=I+\ell=3+2$, the extra equation that forces $W=B+1$.}

\section{A signed count that cannot vanish}\label{sec:count}
\Needspace{7\baselineskip}
\begin{proposition}\label{prop:forest}
Whenever a pair exists, there is a pair $(r,s)$ for which $Q_s$ is a forest.
\end{proposition}
We prove this by evaluating one finite sum in two ways. Fixing $s$ will cancel every cyclic $Q_s$. Fixing the undirected cycles formed by both states will show that a suitable sum is nonzero.

For a pair, let $D(r,s)$ contain the edges $r(t)\to s(t)$. Every vertex of $X$ has one incoming and one outgoing edge, so $D$ is a union of vertex-disjoint directed cycles covering $X$. There are no loops or doubled undirected edges, because the states differ at every face and every edge belongs to just one black face. The cycles are therefore simple. Lemma~\ref{lem:disk} gives each of them odd length.

To assign a sign that reverses on either kind of cycle, put
\begin{equation}\label{eq:sign}
\begin{split}
 J(r,s)&=\frac13\sum_{t\in A}\delta(r(t),s(t))
       =\#\cw\text{ cycles of }D-\#\ccw\text{ cycles of }D,\\
 \eps(r,s)&=(-1)^{(n+J(r,s))/2}.
\end{split}
\end{equation}
The exponent is an integer: if $D$ has $c$ cycles, their odd lengths sum to $n$, so $c\equiv n\pmod2$, and $J\equiv c\pmod2$. The two operations we will use are different.
\begin{center}
\begin{tabular}{@{}lll@{}}\toprule
Cycle reversed & Data changed & Data fixed\\\midrule
A cycle of $Q_s$ & $r$ only & $s$, hence any function of $s$\\
A cycle of $D$ & exchange $r$ and $s$ there & undirected $D$\\\bottomrule
\end{tabular}
\end{center}

First fix $s$ with cyclic $Q_s$ and a partner. Choose a simple cycle $K$ by any fixed ordering of the cycles of $Q_s$, using $s$ alone. The orientation supplied by $r$ makes $K$ directed. On its supplying faces replace $r(t)$ by $u(t)$, the third corner. This cyclically permutes the assigned vertices, so it gives another state $r'$ disjoint from $s$. The move reverses $K$, and applying it twice restores $r$. It has no fixed point. On each changed triangle the cyclic order of its three corners gives
\[
 \delta(u(t),s(t))-\delta(r(t),s(t))=2\delta(r(t),u(t)).
\]
For example, when $(r,u,s)$ is counterclockwise the left side is $1-(-1)=2$; the other cyclic order negates all three signs. Lemma~\ref{lem:disk} on $K$ therefore gives $J(r',s)-J(r,s)=\pm2$, and $\eps(r',s)=-\eps(r,s)$. It follows that for every function $F$ of the second state,
\begin{equation}\label{eq:cancel}
 \sum_{(r,s)\text{ pair}}\eps(r,s)F(s)
 =\sum_{(r,s):\,Q_s\text{ forest}}\eps(r,s)F(s).
\end{equation}
States without a partner contribute no terms. The cycle selection stays the same on the return step because $s$ stays fixed.

In the running example $K$ is $8,1,9,8$. Reversal changes the entries of $r$ at $a_7,a_9,a_{21}$ from $(8,1,9)$ to $(1,9,8)$. All entries of $s$ remain fixed. The values of $J$ are $0$ and $-2$; since $n=10$, the signs are $-1$ and $+1$. These two terms cancel for every $F(s)$.

The unweighted sum cannot prove existence: reversing just one cycle of $D$ also changes the sign, so that sum is zero. We need a function of $s$ that distinguishes the orientations of $D$. A face inside a cycle gives the distinction.

For each white face $w$, choose a point in its interior and a directed polygonal path $\gamma_w$ from that point to a point in $t_0$. Choose the paths to avoid all vertices and black centers and to cross spokes transversely, finitely often. Such paths exist by a small perturbation of polygonal paths; they need not be mutually disjoint. For a spoke from $c_t$ to $v$, count $+1$ when it crosses $\gamma_w$ from its right side to its left side, and $-1$ for the opposite crossing, with sides taken relative to the direction of $\gamma_w$. Call the total $\kappa_w(t,v)$. Define the \emph{height} of $s$ at $w$ by
\begin{equation}\label{eq:height}
 h_s(w)=\sum_{t\in A}\kappa_w(t,s(t)).
\end{equation}
The paths stay fixed. Heights may be negative and depend on the paths; only their changes will matter.

\begin{figure}[htbp]\centering
\begin{tikzpicture}[>=Stealth,font=\footnotesize,scale=1]
\draw[thick,postaction={decorate},decoration={markings,mark=at position .12 with {\arrow{>}}}] (0,0) circle (1.2);
\fill (0,0) circle (1.5pt) node[below left] {$w$};
\draw[->,very thick] (0,0)--(2,0) node[right] {$t_0$};
\node at (.6,.25) {$\gamma_w$};
\node at (1.4,.4) {$+1$};
\node at (0,-1.6) {start inside};
\begin{scope}[xshift=5.8cm]
\draw[thick,postaction={decorate},decoration={markings,mark=at position .12 with {\arrow{>}}}] (0,0) circle (1.2);
\fill (-2,0) circle (1.5pt) node[below] {$w$};
\draw[->,very thick] (-2,0)--(2,0) node[right] {$t_0$};
\node at (-1.4,.4) {$-1$};\node at (1.4,.4) {$+1$};
\node at (0,-1.6) {start outside};
\end{scope}
\end{tikzpicture}

\caption{The height change detects enclosure. A counterclockwise curve crosses an outward path with signed total $+1$ if its starting point lies inside, and $0$ if it lies outside. Extra entrance--exit pairs cancel. Replacing an edge of $D$ by its two-spoke route changes no white face's side of the curve.}\label{fig:height}
\end{figure}

Fix an undirected edge set $D$ arising from a pair. It determines the chosen edge in every black face, since that edge has a unique black incident face. Its pairs are exactly the independent orientations of its cycles. Indeed, in any such orientation tails and heads each use every nonroot once. Let $s_-$ be the second state when all cycles are clockwise, and let $U$ be the set of cycles we reverse to counterclockwise. For a cycle $C$, let $\In(C)$ be the white faces in its bounded disk. Then
\begin{equation}\label{eq:heightchange}
 h_{s_U}(w)=h_{s_-}(w)+\sum_{C\in U}{\bf1}_{w\in\In(C)}.
\end{equation}
To prove this, consider a counterclockwise edge $v\to v'$ of $C$ supplied by $t$. Changing its orientation from clockwise to counterclockwise changes the height contribution by $\kappa_w(t,v')-\kappa_w(t,v)$. This is the signed crossing count of the route $v\to c_t\to v'$. Together these routes form a simple closed curve: the original cycle vertices are distinct, the supplying black centers are distinct, and different spoke interiors are disjoint.

Its orientation and enclosed white faces agree with $C$. Slide each edge through the sector bounded by that edge and its two spokes. The sector excludes the third corner, even if that corner occurs elsewhere on $C$, and its interior meets no other part of the curve. Each slide preserves simplicity and orientation, and passes through no white face. Finally, a counterclockwise simple curve has signed crossing count $+1$ with a path starting inside and ending outside, and zero when both endpoints are outside: the crossings alternate between exit and entrance. With our convention an exit has sign $+1$. This proves~\eqref{eq:heightchange}, including for nested cycles.

We can now choose $F$. Let $m$ be the smallest number of cycles in any $D$ arising from a pair. It exists because there is a pair and the set of pairs is finite; also $m\geq1$. Choose a $D_*$ with $m$ cycles, and choose a white face $w_i$ inside its $i$th cycle for $1\leq i\leq m$. Lemma~\ref{lem:disk} ensures that each cycle contains a white face. Repetitions among the $w_i$ are allowed. Put
\begin{equation}\label{eq:test}
 F(s)=\prod_{i=1}^m h_s(w_i).
\end{equation}
We claim the left side of~\eqref{eq:cancel} is nonzero.

Here is the full counting argument, for arbitrary $m$. Fix an undirected $D$ with cycle set $\mathcal C$, of size $c\geq m$. Write $b_i=h_{s_-}(w_i)$ and $a_{iC}={\bf1}_{w_i\in\In(C)}$. With $U$ its counterclockwise cycles, $J=c-2|U|$. Thus its group contributes
\begin{equation}\label{eq:group}
 (-1)^{(n+c)/2}\sum_{U\subseteq\mathcal C}(-1)^{|U|}
       \prod_{i=1}^m\left(b_i+\sum_{C\in U}a_{iC}\right).
\end{equation}
Expand the product. A term chooses a set $I\subseteq\{1,\ldots,m\}$ of nonconstant factors and, for each $i\in I$, a cycle $\phi(i)$. It occurs precisely for those $U$ containing the image of $\phi$. Summing its signs gives
\[
 \sum_{U\supseteq\operatorname{im}\phi}(-1)^{|U|}
 =(-1)^{|\operatorname{im}\phi|}(1-1)^{c-|\operatorname{im}\phi|}.
\]
This is zero unless $\phi$ is onto. If $c>m$ that is impossible. If $c=m$, it requires $I=\{1,\ldots,m\}$ and $\phi$ a bijection. Therefore~\eqref{eq:group} is zero for $c>m$, and for $c=m$ it is
\begin{equation}\label{eq:positive}
 (-1)^{(n-m)/2}
 \#\{\text{bijections }\beta:\{1,\ldots,m\}\to\mathcal C:
                         w_i\in\In(\beta(i))\text{ for every }i\}.
\end{equation}
Indeed the sign is $(-1)^{(n+m)/2+m}=(-1)^{(n-m)/2}$. All groups now have the same sign or contribute zero, and the chosen $D_*$ contributes at least one, by our choice of the $w_i$. The sum is nonzero. Equation~\eqref{eq:cancel} forces a forest pair to remain, proving Proposition~\ref{prop:forest}.

\begin{example}[What the count sees]
On the cube, $D$ has one triangular cycle enclosing $y_3$. Reversing it increases $h_s(y_3)$ by one, so the two opposite signs give a nonzero difference.
In the larger example there are eight pairs, forming two groups of four orientations of two cycles each. Six pairs have forest $Q_s$ and two cancel as described above. Thus $m=2$. For the displayed $D$, choose $w_1=239$ inside the long cycle and $w_2=78\,12$ inside the small cycle. The second face is inside both cycles. If their height increments are indexed by $u,v\in\{0,1\}$, the four-term difference is
\[
 \sum_{u,v=0}^1(-1)^{u+v}(b_1+u)(b_2+u+v)=1.
\]
This displays why two heights are needed here and why nesting causes no problem. The other group contributes a nonnegative integer with the same overall sign. The proof does not assert that a preferred orientation within the displayed group survives cancellation.
\end{example}
\section{From a forest state to the cycle}\label{sec:recover}
Assume for now that $T$ has no separating triangle. Lemma~\ref{lem:pair} and Proposition~\ref{prop:forest} give a pair $(r,s)$ with forest $Q_s$. By Lemma~\ref{lem:components}, its three components have one root each.

We follow OpenAI's recovery \cite[Section 6]{openai2026}. We still need one selected side at every \emph{white} triangle as well as every black triangle. A tree on the white faces will assign those sides. Let $H$ be the sphere graph of all black centers and their spokes, including the center of $t_0$. It is connected: each edge of $T$ can be routed through its black center. Its faces correspond exactly to the white triangles. In fact the spokes split a black triangle into three sectors adjoining its sides; erasing the edges of $T$ joins each white triangle to those three neighboring sectors, and does not merge different white triangles.

Delete the spokes $(t,s(t))$, $t\in A$. What remains is $Q_s$ with one subdivision vertex on each edge, joined by the three outer spokes to $c_{t_0}$. The roots lie in distinct components, so this is a spanning tree of $H$. By Lemma~\ref{lem:duality}, the duals of the deleted spokes form a tree $W_s$ on the white faces.

Root $W_s$ at the white face $y_0$ across $e_0$ from $t_0$. Assign each tree edge to its endpoint farther from $y_0$. Every white face except $y_0$ receives exactly one edge. A deleted spoke $(t,s(t))$ borders the white faces across the two sides of $t$ incident with $s(t)$. Select the side bordering its assigned white face. Finally select $e_0$, and call the selected set $P$.

Each bounded black face contributes one side, and $t_0$ contributes $e_0$. Every white face except $y_0$ receives its unique parent edge in $W_s$, and $y_0$ receives $e_0$. Thus $P$ has exactly one side of every face. This proves the degree condition for the complementary dual. We must still prove connectivity.
\fig{12-tree-cotree}{1}{The larger example, now using the state $r$ of~\eqref{eq:running} as the forest second state. The dotted spokes are deleted. The remaining spokes form one tree, joined through the outer center. The green edges crossing the deleted spokes form the tree on white faces, rooted at the face across $04$.}

\Needspace{8\baselineskip}
\begin{lemma}\label{lem:selected}
The selected set $P$ is a forest with two components, and those components are induced trees with the prescribed pattern at $t_0$.
\end{lemma}
\begin{proof}
Put $P_0=P\setminus\{e_0\}$ and direct its edge in $t$ out of $s(t)$. Nonroots have outdegree one and roots zero. A cycle in $P_0$ would therefore be directed. For its disk, Lemma~\ref{lem:disk} gives $W-B=1$ and $h-2\ell=3$.

But the selected edges pair black faces with white faces. A face in the disk is paired to another face in the disk unless its selected edge lies on the boundary cycle. Every boundary edge is selected, with $\ell$ black and $h-\ell$ white faces on its inner side. Removing these boundary-paired faces leaves equal numbers of each color, so $B-\ell=W-(h-\ell)$, giving $W-B=h-2\ell=3$, a contradiction. Internal selected chords have both incident faces inside and cause no exception. Thus $P_0$ is a forest.

The outdegree count of Lemma~\ref{lem:components} gives one root in each component of $P_0$. Adding $e_0$ joins two root components and leaves the third apart. Hence $P$ is a forest with two components.

Its complementary dual is connected by Lemma~\ref{lem:duality} and has degree two at every face, so it is one spanning cycle. Draw that cycle crossing each unselected edge of $T$ exactly once and crossing no selected edge. The Jordan curve labels each vertex of $T$ by its side. An edge has same-side endpoints exactly when it is selected: it crosses the curve zero times if selected and once otherwise. Both sides contain vertices, as any unselected edge shows. Each component of $P$ lies on one side, and there are two components, so they are exactly the two side classes. Every edge within a class belongs to $P$. The components are induced trees, with the endpoints of $e_0$ together and the third root apart.
\end{proof}
\fig{13-selected}{1}{The final selection in the larger example. A white face other than the root selects the side corresponding to its parent edge in the white tree; the root selects $04$. The bold edges form two induced trees. The curve crosses every other edge once and visits all 22 faces, giving the Hamiltonian cycle in $G$. The separate degree and connectivity arguments are both needed.}

\begin{example}[The cube completed]
For $s=(5,3,4)$ the white tree has edges $y_0y_3,y_1y_3,y_2y_3$, rooted at $y_0$. Its assignment is
\begin{center}
\begin{tabular}{@{}ccccc@{}}\toprule
Black face & Deleted corner & White-tree edge & Child & Selected side\\\midrule
$t_1=135$ & $5$ & $y_0y_3$ & $y_3$ & $35$\\
$t_2=234$ & $3$ & $y_1y_3$ & $y_1$ & $23$\\
$t_3=045$ & $4$ & $y_2y_3$ & $y_2$ & $04$\\\bottomrule
\end{tabular}
\end{center}
Adding $01$ gives the induced trees on $\{0,1,4\}$ and $\{2,3,5\}$. The complementary cube cycle, naming its vertices by the faces of the octahedron, is
\[
 t_0,y_2,t_2,y_3,t_3,y_0,t_1,y_1,t_0.
\]
It visits all eight vertices and avoids the prescribed edge $t_0y_0$.
The same recovery applies to the larger example using $r$ as the forest second state; Figure~\ref{fig:13-selected} displays the two classes and their complementary cycle.
\end{example}

\section{Gluing across a triangle}\label{sec:glue}
We finish Proposition~\ref{prop:dual} by induction on $|V(T)|$. The construction just completed proves the case with no separating triangle. Otherwise cut along a separating 3-cycle and cap each side by a triangular face. Both capped pieces are simple sphere triangulations with at least four vertices, and both have fewer vertices than $T$.

The inherited face coloring extends to each cap. On one side of the separating triangle, signed side counting makes $3(B-W)$ a sum of three signs $\pm1$. Divisibility by three forces the signs to agree. Thus the three faces incident with the boundary on that side have one color; give the cap the opposite color.

Apply induction first to the piece containing the original face $t_0$, with its prescribed edge $e_0$. Both tree labels occur on the cap, because a monochromatic triangle would be a cycle in an induced tree. Its unique same-label edge determines the prescription on the other piece. Apply induction there with the cap as distinguished face and that edge as the prescribed edge. Rename the two labels to agree on the common triangle.

For each label the two trees meet in a connected subtree, either a single vertex or a single edge with its endpoints. Their union is a tree: the union is connected and its edge count is one less than its vertex count, subtracting the shared subtree once. It is induced because no edge joins the strict interiors of opposite sides of the separating triangle. The original prescription survives. This completes the induction and proves Proposition~\ref{prop:dual}.
\fig{14-gluing}{1}{Gluing two octahedral pieces along $\{3,4,5\}$. The first piece puts $3$ and $5$ together. Prescribing $35$ in the second piece makes its labels agree. One pair of trees intersects in the edge $35$ and the other in the vertex $4$, so both unions are trees.}

\begin{proof}[Proof of Theorem~\ref{thm:main}]
Form the simple face-colored dual triangulation as in Section~1. Choose $t_0$ incident with the dual edge $e_0=e^*$, and apply Proposition~\ref{prop:dual}. The edges within the two tree classes form a spanning forest with exactly one side of every triangle. Their complementary dual is connected by Lemma~\ref{lem:duality} and has degree two at every vertex of $G$, hence is a Hamiltonian cycle. It omits $e$ because $e_0$ lies within one tree class.
\end{proof}

\bibliographystyle{amsplain}
\bibliography{refs}
\end{document}
